% ============================================================
% DRAFT NOTE (target)
% Contains: A plain-language statement of the hazard target, then the formal
% definition: paired same-prefix evaluation, the future-sum divergence family,
% horizons, and the uncompressed-cache convention.
% Written now: The target definition per the 2026-09-01 pivot.
% Pending: Exact target variant chosen as primary (sum vs max, JS vs KL) and
% the frozen horizon set, pending first modeling results.
% Sources: pivot discussion; Jocana/herald-logits dataset card
% ============================================================
\section{The Forecasting Target: Near-Term Compression Damage}
\label{sec:target}

Consider a model decoding with a compressed KV cache. At each step it emits a
token from a next-token distribution that the compressed cache may already
have distorted. The quantity we want online is the near-term hazard: how much
distortion, attributable to compression alone, the generation will accumulate
over its next few tokens. Defining this requires separating compression
damage from ordinary autoregressive divergence, since two runs of the same
model can drift apart for reasons that have nothing to do with the cache.

\paragraph{Paired same-prefix evaluation.}
Let a compressed run under compressor \compressor{} at ratio \ratio{} emit
the realized prefix $y_{1:t}$, and let $p^{\mathrm{c}}_{i}$ denote its
next-token distribution at step $i$. The reference $p^{\mathrm{u}}_{i}$ is
the \emph{uncompressed} model's next-token distribution evaluated on the same
realized prefix $y_{1:i-1}$, by teacher forcing. Because both distributions
condition on identical tokens, their divergence at step $i$ isolates the
effect of the compressed cache: it asks what the full-cache model would have
predicted right here.

\paragraph{The target family.}
The primary target is cumulative near-term divergence at horizon $k$,
\begin{equation}
\fsj{k}(t) \;=\; \sum_{i=t+1}^{t+k}
\mathrm{JS}\!\left(p^{\mathrm{c}}_{i},\, p^{\mathrm{u}}_{i}\right),
\label{eq:target}
\end{equation}
with Jensen--Shannon divergence chosen for boundedness and symmetry.
Companion targets are the KL analogue and the worst single-step divergence
$\max_{i \in (t, t+k]} \mathrm{JS}(p^{\mathrm{c}}_{i}, p^{\mathrm{u}}_{i})$,
and horizons $k \in \{5, 10, 25, 50\}$ span immediate alarm to
paragraph-scale drift. \todo{Freeze the primary variant and horizon set once
first modeling results indicate which the downstream use case needs.}

\paragraph{Defined at every step.}
The target is well-defined whether or not the cache is currently compressed.
For an uncompressed cache the two distributions coincide and the target is
identically zero; such steps are represented with an inactive compressor and
ratio zero. This convention anchors the low end of the scale and lets one
forecaster run continuously across compressed and uncompressed operation. The
target is an oracle quantity: computing it requires the paired uncompressed
evaluation, which is available at dataset-construction time but never at
inference time. It and every future-looking field are training targets only,
never model inputs.

\paragraph{Semantic boundary.}
This is a current-state, actual-damage target: it quantifies the damage the
\emph{already active} compression configuration will inflict over the
forecast window, under fixed compression throughout that window. It is not a
prospective target; forecasting what \emph{would} happen if compression were
activated on a currently uncompressed cache is a separate counterfactual
problem, out of scope here.

\paragraph{Relation to task quality.}
Divergence is measured in distribution space, not task space. We treat final
task quality and catastrophic-failure indicators as secondary validation:
Section~\ref{sec:setup} tests whether predicted distributional damage
corresponds to realized quality failures, rather than assuming it.
